\subsection{Descent data}

Let X and Y be topological spaces, and let \(f \colon \mathrm{X} \to \mathrm{Y}\) be a continuous map. Let \((\mathbf{Z}, p)\) be an X-space.

DEFINITION 1. --- A descent datum relative to f on the X-space (Z, p) is a continuous map \(\tau\colon \mathrm{Z}\times_{\mathrm{Y}}\mathrm{X}\to\mathrm{Z}\) satisfying the following two properties:

\begin{enumerate}
\def\labelenumi{(\roman{enumi})}
\item
  For every pair \((x, x') \in \mathrm{X} \times_{\mathrm{Y}} \mathrm{X}\) , the map \(z \mapsto \tau(z, x')\) induced by restriction is a bijection \(\tau_{x,x'}\) from \(Z_x\) onto \(Z_{x'}\) ;
\item
  For every triple \((x,x',x'')\) of points of X such that \(f(x)=f(x')=f(x'')\) , we have
\end{enumerate}

\[
\tau_ {x, x ^ {\prime \prime}} = \tau_ {x ^ {\prime}, x ^ {\prime \prime}} \circ \tau_ {x, x ^ {\prime}}.
\]

If \(\tau\) is a descent datum on \((\mathrm{Z},p)\) , the family \((\tau_{x,x'})\) , for \((x,x')\in \mathrm{X}\times_{\mathrm{Y}}\mathrm{X}\) is thus a (right) action of the groupoid \(\mathrm{X}\times_{\mathrm{Y}}\mathrm{X}\) on the X-space \((\mathrm{Z},p)\) (II, p.~167). In particular, we have \(\tau_{x,x} = \mathrm{Id}_{\mathrm{Z}_x}\) for all \(x\in \mathrm{X}\) , which is also written \(\tau (z,p(z)) = z\) for all \(z\in \mathrm{Z}\) . Conversely, given a (right) action of the groupoid \(\mathrm{X}\times_{\mathrm{Y}}\mathrm{X}\) on \((\mathrm{Z},p)\) , the map from \(\mathrm{Z}\times_{\mathrm{Y}}\mathrm{X}\) into Z defined by \((z,x)\mapsto z\cdot (p(z),x)\) satisfies the relations (i) and (ii) of the definition.

Let X and Y be topological spaces, \(f\colon X\to Y\) a continuous map and let \((Z,p)\) be an X-space. If \(\tau\) is a descent datum relative to \(f\) on \((Z,p)\) , let \(\mathrm{R}_{\tau}\{z_{1},z_{2}\}\) be the relation defined by \(f(p(z_{1})) = f(p(z_{2}))\) and \(\tau (z_1,p(z_2)) = z_2\) . It is the equivalence relation in Z deduced from the action of the groupoid \(X\times_Y X\) defined by \(\tau\) ; it is compatible with the map \(f\circ p\) . We say that this is the equivalence relation associated with the descent datum \(\tau\) . The canonical bijection \((z_{1},z_{2})\mapsto (z_{1},p(z_{2}))\) of the graph \(\Gamma\) of the equivalence relation \(\mathrm{R}_{\tau}\) onto \(Z\times_Y X\) is a homeomorphism, whose inverse maps an element \((z_{1},x_{2})\in Z\times_Y X\) onto \((z_{1},\tau (z_{1},x_{2}))\) .

Conversely, let R be an equivalence relation in Z compatible with the map \(f \circ p: Z \to Y\) and let \(\Gamma\) be the graph of R. Suppose moreover that the map \(p_{2}: (z_{1}, z_{2}) \mapsto (z_{1}, p(z_{2}))\) defines a homeomorphism from \(\Gamma\) onto \(Z \times_{Y} X\) . The map \(\tau: Z \times_{Y} X \to Z\) given by \(pr_{2} \circ p_{2}^{-1}\) is continuous; it is a descent datum relative to f on \((Z, p)\) and the relation R is the equivalence relation associated with \(\tau\) .

Examples. --- 1) Let X and Y be topological spaces, let \(f\colon X \to Y\) be a continuous map and let \((T,q)\) be a Y-space. Set \(Z = X \times_Y T\) . The map \(\tau\) from \(Z \times_Y X\) into Z which sends \(((x,t),x')\) to \((x',t)\), is a descent datum relative to \(f\) on the X-space \((X \times_Y T, pr_1)\) , called the canonical descent datum. For \(z_1 = (x_1, t_1)\) and \(z_2 = (x_2, t_2) \in Z\) , the relation \(R_{\tau}\{z_1, z_2\}\) is equivalent to \(t_1 = t_2\) .

\begin{enumerate}
\def\labelenumi{\arabic{enumi})}
\setcounter{enumi}{1}
\tightlist
\item
  Let Y be a topological space, \((\mathrm{V}_{i})_{i\in\mathrm{I}}\) a family of subsets of Y, and for every \(i\in I\) , let \((\mathrm{Z}_{i},p_{i})\) be a \(V_{i}\)-space. Denote by X the topological sum of the family \((\mathrm{V}_{i})_{i\in\mathrm{I}}\) and \((\mathrm{Z},p)\) the X-space sum of the family \((\mathrm{Z}_{i},p_{i})\) . Let \(f:X\to Y\) be the canonical map.
\end{enumerate}

The space \(X \times_{Y} X\) is then identified with the sum space of the family \((\mathrm{V}_{i} \cap \mathrm{V}_{j})_{(i,j) \in \mathrm{I} \times \mathrm{I}}(\mathrm{I}, \text{p.~4, example 5})\) . Let \(\tau\) be a descent datum relative to f on \((\mathrm{Z}, p)\) . For every pair \((i, j) \in \mathrm{I} \times \mathrm{I}\) one defines a continuous map \(\tau_{i,j}\colon p_{i}^{-1}(\mathrm{V}_{i}\cap\mathrm{V}_{j})\to p_{j}^{-1}(\mathrm{V}_{i}\cap\mathrm{V}_{j})\;\text{by}\;z\mapsto\tau(z,(p_{i}(z),j))\) . The family \((\tau_{i,j})\) satisfies the following properties:

\begin{enumerate}
\def\labelenumi{(\roman{enumi})}
\item
  For every \(i \in \mathrm{I}\) , we have \(\tau_{i,i} = \mathrm{Id}_{\mathrm{Z}_i}\) ;
\item
  For every pair \((i,j)\in \mathrm{I}\times \mathrm{I},\tau_{i,j}\) is an isomorphism of \((\mathrm{V}_i\cap \mathrm{V}_j)\)-spaces;
\item
  For every triple \((i,j,k)\in\mathrm{I}\times \mathrm{I}\times \mathrm{I}\) and every \(z\in p_i^{-1} (\mathrm{V}_i\cap \mathrm{V}_j\cap \mathrm{V}_k)\) , we have \(\tau_{j,k}(\tau_{i,j}(z)) = \tau_{i,k}(z)\) .
\end{enumerate}

Conversely, any family \((\tau_{i,j})\) possessing the above properties comes from a unique descent datum relative to f on \((Z,p)\) .

